% ---------------------------------------------------------------------------
% 2026-1 기계학습기초 팀 프로젝트 Final Report
% arXiv preprint style (XeLaTeX)
% ---------------------------------------------------------------------------
\documentclass[10pt,a4paper,twocolumn]{article}

% ---------------------------------------------------------------------------
% Engine check
% ---------------------------------------------------------------------------
\usepackage{iftex}
\ifPDFTeX
  \errmessage{This template requires XeLaTeX or LuaLaTeX. In Overleaf, set Compiler to XeLaTeX.}
\fi

% ---------------------------------------------------------------------------
% Language & Font (Korean support via kotex + fontspec)
% ---------------------------------------------------------------------------
\usepackage{fontspec}
\usepackage{kotex}

% ---------------------------------------------------------------------------
% Layout
% ---------------------------------------------------------------------------
\usepackage[
  a4paper,
  top=19.1mm, bottom=25.4mm,
  left=17mm,  right=17mm,
  columnsep=8mm
]{geometry}
\usepackage{setspace}
\usepackage{microtype}
\usepackage{ragged2e}

% ---------------------------------------------------------------------------
% Math
% ---------------------------------------------------------------------------
\usepackage{amsmath}
\usepackage{amssymb}

% ---------------------------------------------------------------------------
% Figures & Tables
% ---------------------------------------------------------------------------
\usepackage{graphicx}
\usepackage{booktabs}
\usepackage{array}
\usepackage{tabularx}
\usepackage{multirow}
\usepackage{caption}
\usepackage{subcaption}

% ---------------------------------------------------------------------------
% Lists
% ---------------------------------------------------------------------------
\usepackage{enumitem}

% ---------------------------------------------------------------------------
% Colors
% ---------------------------------------------------------------------------
\usepackage{xcolor}

% ---------------------------------------------------------------------------
% Section headings  (arXiv / NeurIPS style)
% ---------------------------------------------------------------------------
\usepackage{titlesec}

% ---------------------------------------------------------------------------
% Bibliography  (numbered, in citation order — standard arXiv/ML style)
% ---------------------------------------------------------------------------
\usepackage[numbers,compress]{natbib}

% ---------------------------------------------------------------------------
% Hyperref (load last)
% ---------------------------------------------------------------------------
\usepackage[
  colorlinks = true,
  linkcolor  = black,
  citecolor  = black,
  urlcolor   = blue,
  pdftitle   = {2026-1 기계학습기초 팀 프로젝트 Final Report},
  pdfauthor  = {1팀}
]{hyperref}

% ===========================================================================
% Global typography
% ===========================================================================
\setlength{\parindent}{1em}
\setlength{\parskip}{2pt plus 1pt}
\setstretch{1.05}
\raggedbottom

% Float spacing
\setlength{\textfloatsep}{8pt plus 2pt minus 2pt}
\setlength{\floatsep}{6pt plus 2pt minus 2pt}
\setlength{\intextsep}{6pt plus 2pt minus 2pt}
\setlength{\abovedisplayskip}{4pt plus 1pt}
\setlength{\belowdisplayskip}{4pt plus 1pt}
\setlength{\abovedisplayshortskip}{2pt}
\setlength{\belowdisplayshortskip}{2pt}

% Caption style
\captionsetup{font=small, labelfont=bf, labelsep=period, skip=4pt}

% Hyphenation / line break tolerance
\tolerance=1200
\hyphenpenalty=800
\exhyphenpenalty=800

% ===========================================================================
% Section format  (arXiv / NeurIPS style)
% ===========================================================================
% Numbered, uppercase section titles  →  "1. INTRODUCTION"
\titleformat{\section}
  {\normalfont\normalsize\bfseries}
  {\thesection.}{0.5em}{\MakeUppercase}
\titlespacing*{\section}{0pt}{10pt plus 2pt}{4pt plus 1pt}

% Subsection: bold italic  →  "4.1 Dataset"
\titleformat{\subsection}
  {\normalfont\normalsize\bfseries\itshape}
  {\thesubsection}{0.5em}{}
\titlespacing*{\subsection}{0pt}{6pt}{3pt}

% Subsubsection: italic run-in
\titleformat{\subsubsection}[runin]
  {\normalfont\normalsize\itshape}
  {\thesubsubsection}{0.5em}{}[.]
\titlespacing*{\subsubsection}{0pt}{4pt}{0.5em}

% ===========================================================================
% Table helpers
% ===========================================================================
\setlength{\aboverulesep}{0pt}
\setlength{\belowrulesep}{0pt}
\renewcommand{\arraystretch}{1.15}

% ===========================================================================
% List spacing
% ===========================================================================
\setlist[itemize]{noitemsep, topsep=2pt, partopsep=0pt, leftmargin=1.5em}
\setlist[enumerate]{noitemsep, topsep=2pt, partopsep=0pt, leftmargin=1.5em}

% ===========================================================================
% Title block
% ===========================================================================
\title{%
  \large\bfseries 2026-1 기계학습기초 팀 프로젝트 Final Report
}
\author{%
  홍길동\textsuperscript{\dag}
  \quad 김철수\textsuperscript{\dag}
  \quad 이영희\textsuperscript{\dag}
  \\[4pt]
  \small\textsuperscript{\dag}아주대학교 인공지능융합학과
  \\[2pt]
  \small\texttt{\{hong, kim, lee\}@ajou.ac.kr}
}
\date{%
  \small 제출일: 2026년 6월 12일
  \quad|\quad 담당교수: 교수명 입력
}

% ===========================================================================
\begin{document}

% ---------------------------------------------------------------------------
% Full-width header: title · abstract · keywords
% (single-column section before the two-column body)
% ---------------------------------------------------------------------------
\twocolumn[{%
  \maketitle
  \vspace{-4pt}
  \noindent\rule{\linewidth}{0.8pt}
  \vspace{4pt}
  \begin{center}
    \begin{minipage}{0.96\linewidth}
      \small
      \textbf{Abstract.}\enspace
      본 보고서는 팀 프로젝트의 문제 정의, 관련 연구, 방법론, 실험 설정, 주요 결과,
      해석 및 한계를 종합적으로 정리한 최종 보고서이다. 프로젝트의 목적과 기여를
      명확히 제시하고, 재현 가능한 수준으로 실험 과정을 서술하며, 결과에 대한 분석과
      향후 개선 방향까지 포함하여 작성한다. Abstract는 약 150--200 단어로 간결하게
      작성한다.
    \end{minipage}
  \end{center}
  \vspace{4pt}
  \noindent\textbf{Keywords:} 팀 프로젝트, 기계학습, 실험 분석, 결과 해석, Git 협업
  \vspace{4pt}
  \noindent\rule{\linewidth}{0.8pt}
  \vspace{8pt}
}]

% ===========================================================================
\section{Introduction}
% ===========================================================================

이 절에서는 프로젝트가 다루는 문제의 배경과 필요성을 설명한다. 왜 이 문제가 중요한지,
어떤 실제적 또는 학술적 의미가 있는지를 서술하고, 본 프로젝트에서 해결하고자 하는
과제를 명확히 정의한다. 또한 프로젝트의 주요 목표와 핵심 기여를 요약한다.

포함 권장 흐름:
\begin{itemize}
  \item 문제 배경과 동기
  \item 해결하고자 하는 예측 또는 분류 과제 정의
  \item 프로젝트의 주요 목표와 기여 요약
  \item 이하 절의 구성 안내
\end{itemize}

% ===========================================================================
\section{Related Work}
% ===========================================================================

이 절에서는 기존 연구, 유사한 접근 방식, 혹은 이 문제를 이해하기 위해 필요한
배경지식을 정리한다. 학술 논문, 공개 프로젝트, benchmark, baseline 모델 등을 조사하여
본 프로젝트가 어떤 맥락에 위치하는지 설명한다~\cite{breiman2001random, chen2016xgboost}.

필요에 따라 이 절의 제목은 \emph{Background} 또는 \emph{Problem Setting}으로
조정할 수 있다.

포함 권장 항목:
\begin{itemize}
  \item 기존 연구 또는 유사 프로젝트 요약
  \item 대표적인 접근 방식 비교
  \item 본 프로젝트와 기존 방법의 차이점
  \item 참고한 데이터셋 또는 baseline 소개
\end{itemize}

% ===========================================================================
\section{Approach}
% ===========================================================================

이 절에서는 문제를 해결하기 위해 사용한 전체 접근 방법을 설명한다. 전처리 전략,
모델 선택, 학습 절차, 하이퍼파라미터 설정 등을 독자가 이해할 수 있도록 서술한다.

\subsection{Overall Pipeline}

전체 파이프라인 개요를 설명한다. 필요하다면 그림을 추가하여 데이터 흐름 또는
시스템 구조를 나타낼 수 있다.

\subsection{Preprocessing}

결측치 처리, 정규화, feature engineering 방법을 기술한다.

\subsection{Models}

사용한 모델과 알고리즘, 하이퍼파라미터 설정 근거를 설명한다. 최소 두 개 이상의
기계학습 알고리즘을 비교해야 한다~\cite{pedregosa2011scikit}.

포함 권장 항목:
\begin{itemize}
  \item 베이스라인 모델 설명
  \item 비교 모델 설명
  \item 하이퍼파라미터 설정 근거
  \item 사용한 라이브러리 및 구현 환경
\end{itemize}

% ===========================================================================
\section{Experiment}
% ===========================================================================

이 절에서는 실험 설계와 사용한 데이터셋을 구체적으로 설명한다. 동일한 조건에서
여러 모델을 비교할 수 있도록 데이터 분할, 실험 환경, 평가 지표를 명시한다.

\subsection{Dataset}

\begin{table}[h]
\vspace{-4pt}
\centering
\footnotesize
\setlength{\tabcolsep}{3pt}
\caption{데이터셋 정보}
\begin{tabularx}{\columnwidth}{@{}p{15mm}X@{}}
\toprule
항목 & 내용 \\
\midrule
데이터셋 & 예: UCI Adult / Kaggle 공개 데이터 \\
예측 목표 & 예: 소득 수준 분류 \\
표본 수 & 예: 10{,}000 \\
피처 수 & 예: 18 \\
출처 & URL, 날짜, 라이선스 명시 \\
\bottomrule
\end{tabularx}
\vspace{-4pt}
\end{table}

포함 권장 항목:
\begin{itemize}
  \item 데이터셋 이름과 출처 (URL, 다운로드 날짜, 라이선스)~\cite{dua2019uci}
  \item 전체 샘플 수와 주요 변수 수
  \item 예측 대상(target)과 입력 피처(feature)
  \item 결측치, 이상치, 클래스 불균형 등 데이터 특성
\end{itemize}

\subsection{Experimental Setup}

\begin{table}[h]
\vspace{-4pt}
\centering
\footnotesize
\setlength{\tabcolsep}{3pt}
\caption{실험 설정}
\begin{tabularx}{\columnwidth}{@{}p{18mm}X@{}}
\toprule
항목 & 설정 \\
\midrule
데이터 분할 & train 70\% / val 15\% / test 15\% \\
평가지표 & Accuracy, Macro-F1 \\
베이스라인 & Logistic Regression \\
비교 모델 & Random Forest, XGBoost \\
실험 환경 & Python 3.10, scikit-learn 1.4 \\
\bottomrule
\end{tabularx}
\vspace{-4pt}
\end{table}

포함 권장 항목:
\begin{itemize}
  \item train / validation / test 분할 방식
  \item 교차검증 사용 여부
  \item 실험 환경 (CPU/GPU, 라이브러리 버전)
  \item 평가 지표 및 ablation study 설계
\end{itemize}

% ===========================================================================
\section{Results}
% ===========================================================================

이 절에서는 주요 정량적 결과를 표와 그림 중심으로 제시한다. 단순히 수치를 나열하는
데 그치지 않고, 어떤 모델이 어떤 지표에서 우수했는지 명확하게 드러나도록 구성한다.

\begin{table}[h]
\vspace{-4pt}
\centering
\footnotesize
\setlength{\tabcolsep}{4pt}
\caption{모델별 성능 비교}
\begin{tabular}{@{}lccc@{}}
\toprule
모델 & Accuracy & F1-score & 비고 \\
\midrule
Baseline (LR)  & 0.812 & 0.801 & 기본 설정 \\
Random Forest  & 0.846 & 0.838 & feat.\ eng. 적용 \\
XGBoost        & \textbf{0.861} & \textbf{0.854} & 최종 선택 \\
\bottomrule
\end{tabular}
\vspace{-4pt}
\end{table}

포함 권장 항목:
\begin{itemize}
  \item 모델별 성능 비교 표 (예: Table~3)
  \item confusion matrix, ROC curve, loss curve 등 시각화
  \item 주요 결과 요약 및 baseline 대비 개선 여부
\end{itemize}

% ===========================================================================
\section{Analysis / Discussion}
% ===========================================================================

이 절에서는 결과를 해석하고, 왜 그런 결과가 나왔는지 분석한다. 성능이 잘 나온 이유와
잘 나오지 않은 이유를 모두 다루며, 실패 사례와 한계점도 함께 정리한다. 단순히
``성능이 높았다/낮았다''로 끝내지 말고, 가능한 근거를 들어 설명한다.

포함 권장 항목:
\begin{itemize}
  \item 성능 차이에 대한 원인 분석
  \item 중요한 feature 또는 모델 특성 해석
  \item 오분류/오예측 사례 분석
  \item 데이터 품질 또는 실험 설계의 한계
  \item 개선 방향 및 후속 연구 아이디어
\end{itemize}

% ===========================================================================
\section{Conclusion}
% ===========================================================================

이 절에서는 프로젝트 전체를 간단히 요약한다. 어떤 문제를 다뤘고, 어떤 방법을
사용했으며, 어떤 결과와 의미를 얻었는지를 정리한다. 마지막으로 향후 확장 가능성이나
실제 적용 가능성을 짧게 덧붙일 수 있다.

% ===========================================================================
% References  (arXiv numbered style)
% ===========================================================================
\bibliographystyle{abbrvnat}
\bibliography{references}

% ===========================================================================
% Appendix (Optional)
% ===========================================================================
\appendix

\section*{Appendix}

필요한 경우 부록을 추가할 수 있다.

\begin{itemize}
  \item 추가 실험 결과
  \item 하이퍼파라미터 상세표
  \item 데이터 전처리 세부 절차
  \item 모델 구조 그림
  \item 역할 분담 및 GitHub 협업 로그 요약
\end{itemize}

\end{document}


\usepackage{fontspec}
\usepackage{kotex}
\usepackage{geometry}
\usepackage{setspace}
\usepackage{titlesec}
\usepackage{microtype}
\usepackage{ragged2e}
\usepackage{graphicx}
\usepackage{booktabs}
\usepackage{array}
\usepackage{tabularx}
\usepackage{amsmath}
\usepackage{amssymb}
\usepackage{cite}
\usepackage{caption}
\usepackage{hyperref}
\usepackage{enumitem}

% 폰트는 kotex 기본 설정을 사용해 환경 차이로 인한 깨짐을 방지한다.

\geometry{top=14mm,bottom=17mm,left=14mm,right=14mm,columnsep=9mm}
\hypersetup{
  colorlinks=true,
  linkcolor=black,
  urlcolor=blue,
  citecolor=black,
  pdftitle={2026-1 기계학습기초 팀 프로젝트 Report},
  pdfauthor={학생 이름}
}

\setstretch{1.0}
\setlength{\parindent}{0.8em}
\setlength{\parskip}{2pt plus 1pt}

% 하이픈 및 정렬 개선 (Underfull/Overfull hbox 경고 해결)
\tolerance=1200
\hyphenpenalty=800
\exhyphenpenalty=800
\raggedbottom
\setlength{\textfloatsep}{4pt plus 1pt minus 1pt}
\setlength{\floatsep}{4pt plus 1pt minus 1pt}
\setlength{\intextsep}{4pt plus 1pt minus 1pt}
\setlength{\abovedisplayskip}{4pt}
\setlength{\belowdisplayskip}{4pt}
\setlength{\abovedisplayshortskip}{2pt}
\setlength{\belowdisplayshortskip}{2pt}
\captionsetup{font=small,skip=3pt}

% 섹션 제목 여백 조정 (titlesec)
\titlespacing*{\section}{0pt}{6pt}{3pt}
\titleformat{\section}
  {\normalfont\fontsize{10pt}{10pt}\bfseries}
  {\thesection}
  {0.5em}
  {}

% 테이블 여백 최적화
\setlength{\aboverulesep}{0pt}
\setlength{\belowrulesep}{0pt}

% 리스트 항목 간 간격 축소
\setlist[itemize]{noitemsep, topsep=2pt, partopsep=0pt}
\setlist[enumerate]{noitemsep, topsep=2pt, partopsep=0pt}

\title{2026-1 기계학습기초 팀 프로젝트 Report}
\author{}
\date{}

\begin{document}

\twocolumn[
\vspace*{-1.2em}
\begin{center}
{\LARGE\bfseries 2026-1 기계학습기초 팀 프로젝트 Report\par}
\vspace{0.35em}
{\normalsize 1팀: 홍길동(2022XXXXX), 김철수(2022XXXXX), 이영희(2022XXXXX)\par}
{\small 아주대학교 인공지능융합학과 \quad 담당교수: 교수명 입력 \quad 제출일: \today\par}
\end{center}
\vspace{-0.6em}
\begin{abstract}
본 보고서는 팀 프로젝트의 문제 정의, 데이터 출처, 예상 모델링 전략, 평가 계획, 팀 운영 방식을 1페이지 안에 요약하기 위한 템플릿이다. 프로젝트 안내 문서의 핵심 요구사항을 반영하여 문제 정의의 명확성, 데이터 출처와 사용 가능성, 모델 비교 계획, 역할 분담과 Git 협업 계획을 한 번에 확인할 수 있도록 구성하였다.
\end{abstract}
\vspace{-1.0em}
\noindent\textbf{키워드:} 팀 프로젝트, 보고서, 지도학습, 실험 분석, Git 협업
\vspace{0.5em}
]

\section{문제 정의와 목표}
본 팀은 \textbf{본전공과 연결된 지도학습 문제}를 정의하고, 해당 문제가 실제로 어떤 의사결정에 도움을 줄 수 있는지 설명한다. 문제는 분류 또는 회귀 형태로 명확히 적고, 왜 이 문제가 의미 있는지 한두 문장으로 동기를 제시한다. 예를 들어 ``학과 만족도 예측'', ``장비 고장 여부 분류'', ``수요 예측''처럼 입력과 출력이 분명한 형태로 쓰는 것이 좋다.

\section{데이터 계획}
프로젝트 안내 문서에 맞게 데이터는 공개적으로 수집 가능해야 하며, \textbf{출처 URL, 다운로드 날짜, 라이선스}를 명시해야 한다. 또한 전처리 이후 기준으로 충분한 표본 수와 변수 수를 확보할 계획을 적는다.

\begin{table}[h]
\vspace{-6pt}
\centering
\footnotesize
\setlength{\tabcolsep}{3pt}
\caption{\small 데이터 계획 예시}
\begin{tabularx}{\columnwidth}{@{}p{14mm}X@{}}
\toprule
항목 & 내용 \\
\midrule
데이터셋 & 예: UCI Adult / Kaggle 공개 데이터 \\
예측 목표 & 예: 소득 수준 분류 \\
피처 수 & 10개 이상 확보 예정 \\
표본 수 & 1{,}000개 이상 확보 예정 \\
출처 기록 & URL, 날짜, 라이선스 명시 \\
\bottomrule
\end{tabularx}
\vspace{-6pt}
\end{table}

\section{방법과 평가 계획}
최소 두 개 이상의 기계학습 알고리즘을 비교한다는 점을 미리 밝히고, 동일한 데이터 분할과 동일한 평가 지표를 사용하겠다고 적는다. 예를 들어 로지스틱 회귀와 랜덤포레스트를 비교하고, 정확도와 F1 점수를 주요 지표로 사용할 수 있다. 가능하다면 기준선 모델, 예상 전처리, 변수 선택 방식도 한 줄씩 포함한다.

\section{팀 운영과 산출물 계획}
최종 보고서 단계에서는 팀별 역할 분담과 Git 협업 계획을 짧게 제시하는 편이 좋다. 이는 이후 Git 활용과 보고서 구성 평가에도 직접 연결된다.

\begin{itemize}[leftmargin=1.5em]
  \item 홍길동: 데이터 수집 및 전처리 초안
  \item 김철수: 베이스라인 모델 구현 및 실험 관리
  \item 이영희: 결과 정리, Report 초안, GitHub 문서화
  \item 저장소 운영: feature 브랜치 기반 작업 후 PR로 병합
\end{itemize}

\section{기대 효과와 한계}
마지막 단락에서는 이 프로젝트가 어떤 의미를 가지는지, 그리고 예상되는 제약이 무엇인지 간단히 적는다. 예를 들어 데이터 편향, 클래스 불균형, 피처 부족, 외부 타당성 한계 등을 미리 언급하면 보고서의 완성도가 높아진다.

참고문헌 예시: 공개 데이터셋 소개 문서나 관련 선행 자료를 인용할 수 있다 \cite{daudon2018recurrence}.

\bibliographystyle{ieee_fullname}
\bibliography{references}

\end{document}
